\documentclass[11pt,a4paper]{article}
\usepackage[T1]{fontenc}
\usepackage[utf8]{inputenc}
\usepackage{lmodern,microtype,amsmath,amssymb,amsthm,mathtools}
\usepackage[margin=26mm,headheight=14pt]{geometry}
\usepackage{booktabs,array,enumitem,xcolor,listings,fancyhdr,needspace}
\usepackage[colorlinks=true,linkcolor=blue!45!black,citecolor=blue!45!black,urlcolor=blue!45!black]{hyperref}
\hypersetup{pdftitle={Polynomial Exponents in a Quartic Cardon Product: An All-Degree Classification},pdfauthor={Research manuscript prepared for Keio}}
\pagestyle{fancy}\fancyhf{}
\fancyhead[L]{\small Polynomial exponents in a quartic Cardon product}
\fancyhead[R]{\small 14 November 2026}
\fancyfoot[C]{\thepage}
\setlength{\emergencystretch}{2em}
\setlength{\parskip}{3pt}
\allowdisplaybreaks
\newtheorem{theorem}{Theorem}[section]
\newtheorem{lemma}[theorem]{Lemma}
\newtheorem{corollary}[theorem]{Corollary}
\newtheorem{proposition}[theorem]{Proposition}
\theoremstyle{remark}\newtheorem{remark}[theorem]{Remark}
\newcommand{\R}{\mathbb R}\newcommand{\C}{\mathbb C}
\newcommand{\LP}{\mathcal{LP}}
\newcommand{\Ph}{\Phi_p}
\DeclareMathOperator{\Repart}{Re}
\lstset{language=Python,basicstyle=\ttfamily\scriptsize,breaklines=true,
  columns=fullflexible,keepspaces=true,showstringspaces=false,
  frame=single,rulecolor=\color{black!20},keywordstyle=\color{blue!45!black},
  commentstyle=\color{black!55}}
\title{\textbf{Polynomial Exponents in a Quartic Cardon Product}\\[6pt]
\large An All-Degree Classification and a Prior-Art Audit}
\author{Research manuscript prepared for Keio\\\small with ChatGPT-assisted exploration and verification}
\date{14 November 2026}
\begin{document}
\maketitle
\begin{abstract}
Fix $p(w)=w^4+w^2+1$ and write $p(w)=\prod_{j=1}^4(w+\alpha_j)$.
For an arbitrary real polynomial $Q$, we prove that every coefficient in
$\prod_j e^{Q(x+\alpha_jt)}$ is nonnegative for every real $x$ if and only
if $Q$ is concave on $\R$ and has degree at most four. The same condition
is equivalent to the global product inequality
$\prod_j e^{Q(x+\alpha_jt)}\ge e^{4Q(x)}$ for real $x,t$.
The decisive step is an explicit angular estimate excluding every even
degree at least six; it does not assume a degree bound or rely on a
finite computation. An entire-logarithm argument extends the
classification to all zero-free real entire functions of finite order.
Koide's preprint of 7 October 2026 already contains the degree-at-most-four
classification and the counterexample to Csordas--Vishnyakova
Conjecture 3.6 that motivated this investigation. Those results are
credited here. Our additional conclusion removes the polynomial degree
restriction for this fixed auxiliary polynomial. Historical priority of
that extension is not certified. In particular, this paper is not a
claim to have met all six requested criteria for a new resolution of
a previously named open problem.
\end{abstract}

\section{The result and its precise scope}

Let
\begin{equation}\label{eq:p}
 p(w)=w^4+w^2+1=\prod_{j=1}^4(w+\alpha_j),\qquad
 \{\alpha_j\}_{j=1}^4=\{\omega,\bar\omega,-\omega,-\bar\omega\},
 \quad \omega=e^{i\pi/3}.
\end{equation}
For a real entire function $f$, its Cardon product and coefficients are
\begin{equation}\label{eq:cardon}
 \Ph[f](z,t)=\prod_{j=1}^4 f(z+\alpha_jt)
           =\sum_{k=0}^{\infty}A_k[f](z)t^k.
\end{equation}
The series converges for all complex $z,t$. The symmetry of the four
shifts makes the product even in $t$, and conjugation shows that its
coefficients are real on the real axis. Define the condition
\begin{equation}\label{eq:C}
 \mathsf C_p(f):\qquad A_k[f](x)\ge0
       \quad\text{for every }x\in\R\text{ and every integer }k\ge0.
\end{equation}
The auxiliary polynomial in \eqref{eq:p} is fixed throughout the main
theorem. In particular, no quantifier over all auxiliary polynomials is
being substituted for the original fixed-polynomial condition.

\Needspace{16\baselineskip}
\begin{theorem}[All polynomial degrees]\label{thm:main}
For every $Q\in\R[z]$, the following are equivalent:
\begin{enumerate}[label=\textup{(\roman*)},itemsep=3pt]
\item $\mathsf C_p(e^Q)$ holds.
\item $\Ph[e^Q](x,t)\ge e^{4Q(x)}$ for every $x,t\in\R$.
\item $Q''(x)\le0$ for every $x\in\R$, and $Q$ has degree at most four.
\end{enumerate}
Constants, including the zero polynomial, are included in \textup{(iii)}.
If $q_4=[z^4]Q$, then under these equivalent conditions
\begin{equation}\label{eq:coefficients}
 A_{2n}[e^Q](x)=e^{4Q(x)}
 \sum_{j=0}^{\lfloor n/2\rfloor}
 \frac{(-Q''(x))^{n-2j}(-2q_4)^j}{(n-2j)!\,j!},
 \qquad A_{2n+1}[e^Q](x)=0.
\end{equation}
The usual convention $0^0=1$ is used in this coefficient formula.
\end{theorem}

The theorem answers the auxiliary classification problem for every
polynomial exponent, not just a bounded search of exponents. Its proof
also shows why concavity alone is insufficient: high degree creates a
specific direction in which the global product inequality fails.

\paragraph{Relation to the literature.}
Cardon's theorem imposes a genus-zero-or-one hypothesis after removing
a Gaussian factor \cite{Cardon}. Csordas and Vishnyakova proposed
discarding that hypothesis in Conjecture 3.6 \cite[p.~1649]{CV}.
Koide subsequently supplied counterexamples and, for quartic auxiliary
polynomials, classified exponents of degree at most four
\cite[Section~5]{Koide}. For \eqref{eq:p}, that classification is exactly
condition \textup{(iii)} once the degree restriction is assumed.
The additional part proved here is that no higher-degree polynomial
exponent can satisfy \eqref{eq:C}. The prior counterexample is not
claimed as a discovery of this paper.

\section{Power sums and a degree-exclusion estimate}

Put $s_k=\sum_j\alpha_j^k$. Elementary root-of-unity summation gives
\begin{equation}\label{eq:powers}
 s_0=4,\qquad s_{2k+1}=0,\qquad
 s_{2k}=\begin{cases}4,&3\mid k,\\-2,&3\nmid k.\end{cases}
\end{equation}
For example, $s_2=s_4=-2$ and $s_6=4$.
For any polynomial $Q$ we therefore have the exact finite identity
\begin{equation}\label{eq:newton}
 \log\frac{\Ph[e^Q](x,t)}{e^{4Q(x)}}
 =\sum_{k=1}^{\deg Q}\frac{s_k}{k!}Q^{(k)}(x)t^k.
\end{equation}
Here the logarithm denotes the displayed real exponent when $x,t$ are
real, so no logarithm-branch issue is involved. In particular,
\begin{equation}\label{eq:A2}
 A_2[e^Q](x)=-Q''(x)e^{4Q(x)}.
\end{equation}
Identity \eqref{eq:newton} by itself does not bound the degree. The
following estimate supplies that missing step.

\begin{lemma}[A direction detecting every high even degree]\label{lem:ray}
For every even integer $d\ge6$, there is $r\ge0$ such that
\begin{equation}\label{eq:gap}
 \Delta_d(r):=\sum_{j=1}^4(r+\alpha_j)^d-4r^d>0.
\end{equation}
For $d=6$ one may take $r=0$, giving $\Delta_6(0)=4$.
For $d=2n\ge8$, the explicit choice
\begin{equation}\label{eq:r}
 \theta=\frac\pi n,\qquad
 r=\frac{\sin(\pi/3-\theta)}{\sin\theta}
\end{equation}
satisfies the quantitative bound $\Delta_d(r)\ge4r^d>0$.
\end{lemma}

\begin{proof}
The assertion for $d=6$ follows from \eqref{eq:powers}. Now let $n\ge4$.
Then $0<\theta\le\pi/4<\pi/3$, so $r>0$. Since
\[
 r+\tfrac12=\tfrac{\sqrt3}{2}\cot\theta,
\]
the argument of $r+\omega$ is $\theta$. Define
\[
 R_+=(r^2+r+1)^{1/2},\qquad R_-=(r^2-r+1)^{1/2}.
\]
The two conjugate numbers $r+\omega,r+\bar\omega$ have $d$th powers
equal to the positive real number $R_+^d$, because $d\theta=2\pi$.
The other conjugate pair contributes at least $-2R_-^d$. Consequently
\begin{equation}\label{eq:dominance}
 \sum_j(r+\alpha_j)^d\ge2(R_+^d-R_-^d).
\end{equation}
With $a=1+r^{-2}>1$ and $b=r^{-1}>0$, the binomial theorem gives
\begin{align}
 \frac{R_+^d-R_-^d}{r^d}
  &=(a+b)^n-(a-b)^n \notag\\
  &=2\sum_{\substack{1\le j\le n\\j\ \mathrm{odd}}}
       \binom nj a^{n-j}b^j
   \ge 2na^{n-1}b
   =\frac{2n}{r}(1+r^{-2})^{n-1}
   \ge\frac{2n}{r}.\label{eq:binomial}
\end{align}
The intermediate expression in \eqref{eq:binomial} is simply the
$j=1$ summand, $2na^{n-1}b$; all other displayed summands are positive.
The elementary concavity bound $\sin u\ge2u/\pi$ for
$0\le u\le\pi/2$ yields
\[
 r\le\frac1{\sin(\pi/n)}\le\frac n2.
\]
Combining this with \eqref{eq:dominance}--\eqref{eq:binomial} proves
\[
 \sum_j(r+\alpha_j)^d\ge\frac{4n}{r}r^d\ge8r^d.
\]
Subtracting $4r^d$ gives the claimed strict gap.
\end{proof}

\section{Proof of the classification}

\begin{proof}[Proof of Theorem~\ref{thm:main}]
Suppose \textup{(i)} holds. The odd coefficients vanish. For real $t$,
each remaining term $A_{2k}(x)t^{2k}$ is nonnegative; convergence of the
entire series therefore gives \textup{(ii)}.

Suppose next that \textup{(ii)} holds. By conjugation, the finite sum
$\sum_jQ(x+\alpha_jt)$ is real when $x,t$ are real. Taking logarithms
of the positive product yields
\begin{equation}\label{eq:global}
 D_Q(x,t):=\sum_jQ(x+\alpha_jt)-4Q(x)\ge0
                   \qquad(x,t\in\R).
\end{equation}
The expansion at $t=0$ is
$D_Q(x,t)=-Q''(x)t^2+O(t^4)$. Dividing by $t^2$ and taking the limit
shows that $Q''(x)\le0$ for all real $x$.

If $Q$ is not affine, a polynomial nonpositive on the whole real axis
has even degree and negative leading coefficient when nonconstant.
Applied to $Q''$, this implies that $\deg Q=d$ is even and the leading
coefficient $q_d$ of $Q$ is negative. This also covers the quadratic
case, where $Q''$ is a negative constant.

Assume for contradiction that $d\ge6$. Choose $r$ from
Lemma~\ref{lem:ray} and substitute $x=Rr,t=R$ into \eqref{eq:global}.
As $R\to+\infty$,
\begin{equation}\label{eq:scaling}
 D_Q(Rr,R)=q_d\Delta_d(r)R^d+O(R^{d-1}).
\end{equation}
The leading coefficient on the right is strictly negative. Hence
$D_Q(Rr,R)<0$ for every sufficiently large $R$, contradicting
\eqref{eq:global}. Thus $d\le4$, proving \textup{(iii)}.

Finally assume \textup{(iii)}. Set $q_4=[z^4]Q$, with $q_4=0$ for
lower degrees. Concavity implies $q_4\le0$. The exact identity
\eqref{eq:newton} now reduces to
\begin{equation}\label{eq:quartic}
 \Ph[e^Q](x,t)=e^{4Q(x)}
       \exp\bigl(-Q''(x)t^2-2q_4t^4\bigr).
\end{equation}
Both coefficients in the last exponent are nonnegative. Expanding its
two exponential factors gives \eqref{eq:coefficients}, hence
\textup{(i)}. This completes all implications.
\end{proof}

\begin{remark}[Where the proof adds information]
Equation \eqref{eq:quartic} is the previously established bounded-degree
calculation. The additional conclusion comes from
\eqref{eq:dominance} and \eqref{eq:scaling}: a positive gap independent
of all lower coefficients excludes every high degree. No finite degree
cutoff, numerical sign check, or conjectural growth estimate is used.
\end{remark}

\section{Zero-free entire functions of arbitrary finite order}

The order of an entire function is understood in its usual growth
sense. In the next argument only the following consequence is needed:
if $f$ has finite order $\rho$, then for every $\tau>\rho$ there is
$C_\tau$ such that
\[
 \log|f(z)|\le C_\tau(1+|z|^\tau)\qquad(z\in\C).
\]

\begin{lemma}[Polynomial logarithm]\label{lem:log}
Every zero-free real entire function $f$ of finite order has the form
$f=\sigma e^Q$, where $\sigma\in\{-1,1\}$ and $Q\in\R[z]$.
\end{lemma}
\begin{proof}
Since $f$ never vanishes on $\R$, its real sign $\sigma$ is constant.
The function $f'/f$ is entire. On the simply connected plane define
\[
 h(z)=\log|f(0)|+\int_0^z\frac{f'(w)}{f(w)}\,dw.
\]
The integral is path independent, $f=\sigma e^h$, and $h$ is real on
$\R$. Thus $\Repart h(z)=\log|f(z)|$ has the growth bound above.
Here is the needed elementary disk form of the classical
Borel--Carath\'eodory estimate. Put $g=h-h(0)$ and
$M=\max_{|z|\le2R}\Repart g(z)$. If $M=0$, the maximum principle
makes $g$ constant. If $M>0$, the analytic function
$g/(2M-g)$ has modulus at most one on the disk and vanishes at zero.
Schwarz's lemma bounds its modulus by $1/2$ on $|z|\le R$; hence
$|g|\le2M$ there. The growth bound on $\Repart h$ therefore gives
\[
 \max_{|z|\le R}|h(z)|=O_\tau(1+R^\tau).
\]
For each integer $k>\rho$, choose $\rho<\tau<k$. Cauchy's estimate
then gives
\[
 \frac{|h^{(k)}(0)|}{k!}=O_\tau(R^{\tau-k})\longrightarrow0.
\]
All such derivatives vanish. Hence $h=Q$ is a polynomial, whose
coefficients are real because $h$ is real on $\R$.
\end{proof}

\Needspace{14\baselineskip}
\begin{corollary}[Finite-order zero-free classification]\label{cor:entire}
Let $f$ be a zero-free real entire function of finite order. The following
are equivalent:
\begin{enumerate}[label=\textup{(\roman*)},itemsep=3pt]
\item $\mathsf C_p(f)$ holds.
\item $\Ph[f](x,t)\ge f(x)^4$ for every real $x,t$.
\item $f=\sigma e^Q$, where $\sigma\in\{-1,1\}$, $Q\in\R[z]$,
$\deg Q\le4$, and $Q''\le0$ on $\R$.
\end{enumerate}
In particular, such a function has order at most four. Order four is
attained, for example, by $e^{-z^4-z^2}$.
\end{corollary}
\begin{proof}
Use Lemma~\ref{lem:log}. Since $\sigma^4=1$, the Cardon product is
unchanged by $\sigma$. Apply Theorem~\ref{thm:main}.
\end{proof}

This corollary concerns zero-free functions. Neither a classification
of functions with arbitrary zero divisors nor a classification of
zero-free functions of infinite order is asserted.

\section{Explicit certificates and the known counterexample}

\begin{corollary}[Finite algebraic decision rule]\label{cor:rule}
Write $Q=\sum_{j=0}^d q_jz^j$. Condition $\mathsf C_p(e^Q)$ holds
exactly when all $q_j$ with $j\ge5$ vanish and one of these alternatives
holds:
\begin{align*}
 &q_4=q_3=0,\qquad q_2\le0;\quad\text{or}\\
 &q_4<0,\qquad 3q_3^2\le8q_4q_2.
\end{align*}
The coefficients $q_0,q_1$ are unrestricted.
\end{corollary}
\begin{proof}
Theorem~\ref{thm:main} reduces the question to the nonpositivity of
$Q''(x)=12q_4x^2+6q_3x+2q_2$. A quadratic with negative leading
coefficient is nonpositive everywhere exactly when its discriminant
is nonpositive. The case $q_4=0$ is the elementary linear case.
\end{proof}

For degree four this rule was already included in Koide's
classification \cite{Koide}. Here the rule applies without an input
degree restriction, because the theorem proves that restriction.
Equivalently, every admissible genuine quartic exponent can be written
\[
 Q(z)=-a(z-h)^4-b(z-h)^2+\ell z+c,
 \qquad a>0,\quad b\ge0,\quad h,\ell,c\in\R.
\]
Indeed completing the square in $Q''$ and integrating twice gives this
normal form.

\paragraph{The prior counterexample, reproduced with attribution.}
For $Q(z)=-z^4-z^2$, equation \eqref{eq:quartic} becomes
\begin{equation}\label{eq:known}
 \Ph[e^{-z^4-z^2}](x,t)
 =e^{-4x^4-4x^2}\exp\bigl((12x^2+2)t^2+2t^4\bigr).
\end{equation}
Consequently
\begin{equation}\label{eq:strict}
 A_{2n}(x)\ge e^{-4x^4-4x^2}\frac{(12x^2+2)^n}{n!}>0
 \qquad(x\in\R,\ n\ge0).
\end{equation}
This exact example and its strict positivity are already in Koide's
Section~2 \cite{Koide}; they are not new results here.

For completeness there is a short independent certificate of failure
of membership in $\LP$. If $f\in\LP$, all coefficients in
$f(x+it)f(x-it)$ are nonnegative for real $x$: this is immediate for
real-rooted polynomials from
$\prod_j((x-r_j)^2+t^2)$ and passes to locally uniform limits by
Cauchy's coefficient formula. For the function in \eqref{eq:known},
\[
 f(it)f(-it)=\exp(2t^2-2t^4),\qquad
 [t^6]\exp(2t^2-2t^4)=\frac{2^3}{3!}-4=-\frac83.
\]
Thus this function passes every quartic Cardon inequality but fails
one classical Laguerre inequality. The implication in the original
2013 conjecture is false, as the prior paper already established.
The counterexample is zero-free: the failure concerns membership in
$\LP$, not the existence of a nonreal zero.

\paragraph{An explicit high-degree failure.}
The concave exponent $Q(z)=-z^6$ does not pass. At $x=0$,
$\Ph[e^{-z^6}](0,t)=e^{-4t^6}$, whose $t^6$ coefficient is negative.
For degrees $8,10,\ldots$, Lemma~\ref{lem:ray} supplies a uniform
alternative to guessing an offending coefficient. It locates a real
pair $(x,t)$ violating the sum inequality; the equivalence then
guarantees a negative Taylor coefficient at that $x$.

\section{Corpus synthesis, attribution, and the six criteria}

The investigation used the supplied Verifiable Mathematical Reasoning
Corpus at commit
\texttt{38348353416f1f7fc50cecb93c3b554a00260c13} \cite{VMRC}.
Three relevant themes were root-power traces in item [075], exponential
coefficient positivity in item [119], and exact converses for real-zero
criteria in item [137]. These suggested looking at finite root sums,
positive exponential coefficients, and converse failures together.
Every mathematical input needed for the present theorem is reproved
above, using the standard maximum principle, Schwarz's lemma, and
Cauchy's estimate. No unreviewed corpus
theorem is imported as a black box.

The sequence of research steps matters for attribution. The strict
counterexample was independently reconstructed during this session.
A subsequent primary-source check found the identical pair and formula
in Koide's earlier public source. That candidate was therefore removed
from the list of new conjecture resolutions. The high-degree exclusion
was then proved separately and compared with the earlier bounded-degree
theorem. Independence of derivation does not erase prior publication.

\begin{center}
\begin{tabular}{@{}p{0.12\textwidth}p{0.82\textwidth}@{}}
\toprule
Criterion & Assessment of this deliverable \\
\midrule
G1 & The original Conjecture 3.6 is a pre-existing named target. The
all-degree classification proved here was not located as a previously
named open problem. \\
G2 & The original definitions were retained in checking the counterexample.
The classification extension is explicitly a different auxiliary result. \\
G3 & The original conjecture was already refuted in the inspected
7 October 2026 preprint. Its refutation cannot be counted anew. \\
G4 & Theorem~\ref{thm:main} and Corollary~\ref{cor:entire} are unconditional
and cover all stated degrees, indices, and real variables. \\
G5 & The degree-exclusion estimate is a substantive additional argument;
the bounded-degree formula alone does not give it. \\
G6 & The proposed counterexample has an exact prior duplicate. No exact
duplicate of the all-degree extension was found in the bounded search,
but this does not certify historical priority. \\
\bottomrule
\end{tabular}
\end{center}

\paragraph{Outcome.}
The simultaneous G1--G6 goal was not achieved. The deliverable is a
complete classification theorem with an explicit extension beyond the
inspected prior theorem, accompanied by a documented rejection of a
duplicate conjecture claim.

\clearpage
\appendix
\section{Exact finite verification program}\label{app:verify}

The following standalone Python program uses only integer and rational
arithmetic. It checks Newton sums, the generic quartic identity on a
linear basis, the positive-coefficient recurrence, the ordinary negative
Laguerre coefficient, and exact sample directions. These finite checks
detect transcription errors; Lemma~\ref{lem:ray}, not this program,
proves the assertion for every degree. The program can be copied from
this single \texttt{.tex} file and run with Python 3.

\begin{lstlisting}
from fractions import Fraction as F
from math import comb, factorial
# Newton sums for roots of w^4+w^2+1.
s = [4, 0, -2, 0, -2]
for k in range(5, 101):
    s.append(-s[k-2] - s[k-4])
assert all(s[k] == (0 if k % 2 else
                   (4 if k % 6 == 0 else -2))
           for k in range(101))
# The exponent identity is linear in Q: check the basis 1,z,...,z^4.
# Keys are (power of x, power of t).
for degree in range(5):
    direct = {(degree-k, k): comb(degree, k)*s[k]
              for k in range(degree+1) if s[k]}
    expected = {(degree, 0): 4}
    if degree >= 2:
        expected[(degree-2, 2)] = -degree*(degree-1)
    if degree == 4:
        expected[(0, 4)] = -2
    assert direct == expected
# exp(u*y+v*y^2): positive explicit sum versus differential recurrence.
comparisons = 0
for x in map(F, [-3, -1, 0, 1, 2]):
    u, v = 12*x*x+2, F(2)
    c = [F(1)]
    for n in range(1, 41):
        c.append((u*c[n-1]+(2*v*c[n-2] if n >= 2 else 0))/n)
    for n in range(41):
        explicit = sum((u**(n-2*j)*v**j /
                        (factorial(n-2*j)*factorial(j))
                        for j in range(n//2+1)), F(0))
        assert c[n] == explicit
        assert c[n] >= u**n/factorial(n) > 0
        comparisons += 1
# Classical Laguerre coefficient at the origin.
laguerre_3 = F(2**3, factorial(3)) + F(2)*F(-2)
assert laguerre_3 == F(-8, 3)
def gap(d, r):
    return sum((F(comb(d, k))*s[k]*r**(d-k)
                for k in range(1, d+1)), F(0))
# Exact witnesses for initial high even degrees.
witnesses = [(6, F(0)), (8, F(1, 2)),
             (10, F(1, 2)), (12, F(1))]
for d, r in witnesses:
    assert gap(d, r) > 0
# Quartic decision examples: (q4,q3,q2,expected).
for q4, q3, q2, answer in [(-1,0,-1,True),
        (-1,4,-6,True), (-1,1,0,False),
        (0,0,-1,True), (0,1,-1,False), (1,0,-1,False)]:
    passed = ((q4 == q3 == 0 and q2 <= 0) or
              (q4 < 0 and 3*q3*q3 <= 8*q4*q2))
    assert passed == answer
print('Newton sums through 100: PASS')
print('Quartic exponent identity: PASS')
print('Exact coefficient comparisons:', comparisons)
print('Classical L_3(0):', laguerre_3)
print('Four exact high-degree witnesses: PASS')
print('Finite algebraic decision examples: PASS')
\end{lstlisting}

\section{Reproducibility and source checks}

The source of this manuscript is one self-contained \texttt{.tex} file,
with no external bibliography, figures, or input files. Compile it twice
with \texttt{pdflatex}. The included verifier is not required for
compilation. The recorded exact checks completed without a failed
assertion. Independent mathematical reviews of the angular estimate
and finite-order extension found no gap; these are internal checks,
not external peer review or formal proof-assistant certification.

The original 2013 PDF was inspected at printed page 1649, including the
definitions preceding Conjecture 3.6 and the different hypotheses of
Theorem 3.5. Koide's public metadata and full \texttt{.tex} source were
inspected: the record is dated 7 October 2026, and Sections 2 and 5
contain the overlapping counterexample and bounded-degree theorem.
Targeted searches included the conjecture number, authors, exact paper
title, Cardon growth conditions, and higher-degree polynomial
exponents. The limitations of this search are recorded in the G6 row.

\begin{thebibliography}{9}
\bibitem{Cardon}
D.~A. Cardon,
\emph{Extended Laguerre inequalities and a criterion for real zeros},
in \emph{Progress in Analysis and its Applications}, World Scientific,
2010, pp.~143--149. Preprint arXiv:0911.1122 (2009).
\url{https://arxiv.org/abs/0911.1122}.
Author-hosted text:
\url{https://mathdept.byu.edu/cardon/papers/criterion_for_reality_arXiv.pdf}.

\bibitem{CV}
G.~Csordas and A.~Vishnyakova,
\emph{The generalized Laguerre inequalities and functions in the
Laguerre--P\'olya class}, Central European Journal of Mathematics
\textbf{11} (2013), 1643--1650.
\url{https://doi.org/10.2478/s11533-013-0269-x}.
Conjecture 3.6 is on printed p.~1649.

\bibitem{Koide}
A.~Koide,
\emph{Counterexamples and Sharp Growth Repairs for Cardon--Laguerre
Inequalities}, public preprint, version 1, 7 October 2026.
\url{https://doi.org/10.6084/m9.figshare.34158516.v1}.
The public source, including Sections 2 and 5, was inspected on
14 November 2026. This citation records a preprint, not a claim of
peer-reviewed publication.

\bibitem{VMRC}
QEEDDD, \emph{Verifiable Mathematical Reasoning Corpus}, research drafts,
commit \texttt{38348353416f1f7fc50cecb93c3b554a00260c13},
retrieved 14 November 2026.
\url{https://github.com/QEEDDD/Verifiable-Mathematical-Reasoning-Corpus}.
Relevant exploratory themes: [075] \emph{Shifted Mills and Cubic Pisot
Traces}; [119] \emph{Mathematical Research Manuscript}; [137]
\emph{A finite-difference measure converse and a Jacobi-symbol
determinant counterexample}. The repository describes its contents as
research drafts, not peer-reviewed or formally verified theorems.
\end{thebibliography}
\end{document}
